\section{Replicating Manifesto Project Expert Scores}\label{sec:min-manifesto-results}

Political manifestos are a long-document measurement problem with a
real-world target attached. A party's election platform runs tens
of thousands of words across many policy areas, and political
scientists have spent decades distilling those documents into a
small set of dimension scores that downstream models of voting,
government formation, and policy outcomes take as inputs.
\citet{BenoitEtAl2025} recently posed the question of whether large
language models can replicate that measurement reliably enough to be
used in place of expert coders, and built a benchmark to test the
question on. We approach that benchmark with the C-tree pipeline of
Section~\ref{sec:min-framework}. Each manifesto is split into
chunks. A learned summarizer compresses adjacent chunks recursively
up to one document-level summary. A learned scorer reads that
summary and outputs a number on each of the six policy dimensions.
The training target is the expert-survey mean, and the held-out
evaluation is the Pearson correlation against those means.

\subsection{Data and Replication Target}\label{sec:min-manifesto-data}

The Manifesto Project (MARPOR / CMP) has coded party manifestos
from parliamentary democracies for nearly four decades
\citep{BudgeRobertsonHearl1987, BudgeEtAl2001, KlingemannEtAl2006,
VolkensEtAl2013}. Trained coders break each manifesto into
\emph{quasi-sentences} (the smallest text spans expressing a single
policy idea) and tag each one with a category from a fixed scheme
covering economic policy, social policy, immigration, European
integration, environment, decentralization, and a few dozen other
topics. The current data release reports $5{,}285$ coded platforms
and $2{,}269{,}000$ annotated quasi-sentences
\citep{ManifestoProject2025a, ManifestoCorpus2025, MerzEtAl2016}.

The benchmark we follow, from \citet{BenoitEtAl2025}, pairs the
coded platforms with expert surveys. For each of the $235$ platforms
in their main benchmark, country specialists are asked to place the
party on a seven-point scale along each of six policy dimensions.
Economic policy runs from high public spending and high taxation
through to low public spending and low taxation. Social policy runs
from liberal to conservative on issues like gender equality and
civil liberties. Immigration runs from open to restrictive,
European integration from pro-EU to anti-EU, environment from
environmental protection prioritized through to economic growth
prioritized, and decentralization from centralized governance
through to regional autonomy. The expert-survey value for a
platform on a dimension is the average of those expert placements,
and that average is the target every method below is scored
against.

\citet{BenoitEtAl2025} pose the LLM replication problem on this
benchmark. Frontier proprietary LLMs come close to expert agreement
on most dimensions when run as an 18-score ensemble across three
models, two shot settings, and three summary styles. Open-weight
single-model replications track that ensemble closely on some
dimensions and noticeably less so on others, with decentralization
the consistent breakdown across proprietary and open-weight LLMs
alike.

We approach the same benchmark with a single-model pipeline. A
learned summarizer $g$ produces a short summary of each chunk of a
manifesto and recursively merges adjacent summaries until a single
root summary represents the whole document; a learned scorer $f$
reads that root summary and returns the seven-point score for each
policy dimension. Both $f$ and $g$ are prompt-and-demonstration
wrappers around the same small quantized open-weight
LLM.\footnote{Gemma-4-31B in NVFP4 (4-bit) quantization.
Re-running the dimension-specific pipeline on Gemma-3-27B-IT (the
model \citet{BenoitEtAl2025} use for their open-weight baseline)
gives Pearson correlations close to the Gemma-4 numbers reported
below, indicating the gains here come from the tree structure, not
from a stronger underlying model.}
\S\ref{sec:min-manifesto-setup} below gives the construction in
detail.

When we train $f$ and $g$ for one policy dimension at a time, this
pipeline meets or beats Benoit's proprietary 18-score ensemble on
every dimension we have completed end-to-end. The two completed so
far are economic policy and decentralization; on economic the
single-dimension pipeline clears not only the ensemble but
split-expert agreement, and on decentralization it still clears
the ensemble by roughly $0.07$.
\fixme{Single-dimension runs for the remaining four dimensions
are in flight.}

A natural follow-up question, taken up in
\S\ref{sec:min-manifesto-combined}, is whether one shared
summarizer can stand in for the six dimension-specific ones: can a
single \emph{universal} $g$ produce one summary per manifesto for
all six per-dimension scorers to read? The joint training run shows
that summaries are coupled to the scorer they are built for. A $g$
optimized against an average of six scorer-rewards aligns with each
individual scorer less well than a $g$ optimized for that scorer
alone, and the alignment gap is largest on decentralization. Across
both training regimes, held-out Pearson correlation is stable to how each
manifesto is chunked, in line with Corollary~\ref{cor:schedule}.

\subsection{Setup}\label{sec:min-manifesto-setup}

Each manifesto is split into contiguous chunks (\emph{leaves}) of a
fixed token size. The summarizer $g$ writes a short summary of each
leaf. Adjacent leaf summaries are merged pairwise by the same $g$
into a parent summary, parent summaries are merged into
grandparents, and so on, until a single \emph{root summary}
represents the whole manifesto. The scorer $f$ reads the root
summary and outputs a number on the seven-point scale for each
policy dimension.

Both $f$ and $g$ are DSPy programs that wrap the same underlying
LLM (Gemma-4-31B-NVFP4) with a prompt and a small bank of bootstrap
few-shot demonstrations. We do not hand-craft any demonstrations or
seed examples. We use DSPy's MIPROv2 optimizer to propose and
select prompt instructions and bootstrap demonstrations against a
training metric; the LLM weights themselves are frozen. When we say
we ``train'' $f$ or $g$, we mean MIPROv2 has searched over prompt
instructions and demonstrations and kept the candidate that
maximizes the metric on a held-back validation slice.

Prompt-and-demonstration optimization is one choice among several.
Rather than searching over prompt instructions, we could adapt the
LLM weights directly with LoRA, full fine-tuning, or RL-style
policy optimization against the same metric; the framework's
analysis applies to any optimizer that takes $f$ and $g$ as
optimizable objects. We work with prompts and demonstrations here
because the model weights stay frozen on the inference server and
prompt-level updates are sufficient to make the comparisons we
report.

Training proceeds in alternating rounds, and we use the shorthand
$f^a g^b$ for the $(f, g)$ pair after $a$ scorer-side updates and
$b$ summarizer-side updates. The starting $f^0$ and $g^0$ are DSPy
programs with their default prompts and no optimized
demonstrations. $f^1 g^0$ is the result of training $f$ against
the initial $g^0$, with $g^0$ held fixed; $f^1 g^1$ is the result
of training $g$ against $f^1$, with $f^1$ held fixed; and so on.
Round 1 ends at $f^1 g^1$, round 2 at $f^2 g^2$, etc. Below we
report two points from each leaf's run: the state after round 1
and the best round at that leaf size.

We divide the Benoit benchmark documents into training,
validation, and test sets at the document level: each manifesto
appears in only one of the three. Document counts vary by
dimension because not every platform has an expert score on every
dimension; Table~\ref{tab:min-manifesto-data-summary} gives the
per-dimension breakdown together with the score distribution on
the test set, the proprietary 18-score ensemble's published
Pearson correlation \citep{BenoitEtAl2025}, and our best Pearson
correlation on the same held-out documents. The two
single-dimension runs (economic and decentralization, used for
comparison in \S\ref{sec:min-manifesto-combined}) draw from the
same training pool with only that dimension's expert scores used
as the document target. All numbers below are on the test split.

\begin{table}[!htb]
  \centering
  \begin{tabular}{lrrrcccc}
    \toprule
    & \multicolumn{3}{c}{Documents} & \multicolumn{2}{c}{Test scores (1--7)} & \multicolumn{2}{c}{Best Pearson} \\
    \cmidrule(lr){2-4} \cmidrule(lr){5-6} \cmidrule(lr){7-8}
    Dimension & Train & Val & Test & Mean & SD & Benoit & C-TreePO \\
    \midrule
    Economic         & $140$ & $30$ & $48$ & $4.91$ & $2.45$ & $0.870$ & $0.888$ \\
    Social           & $140$ & $30$ & $48$ & $4.01$ & $2.34$ & $0.920$ & $0.886$ \\
    Immigration      & $125$ & $27$ & $38$ & $5.03$ & $2.47$ & $0.890$ & $0.938$ \\
    EU integration   & $121$ & $26$ & $38$ & $5.09$ & $1.92$ & $0.910$ & $0.965$ \\
    Environment      & $117$ & $25$ & $44$ & $5.10$ & $1.90$ & $0.820$ & $0.795$ \\
    Decentralization & $140$ & $30$ & $48$ & $4.58$ & $1.47$ & $0.490$ & $0.557$ \\
    \bottomrule
  \end{tabular}
  \caption{\textbf{Per-dimension benchmark and headline.} Document
  counts in each split with an expert score on the dimension; mean
  and standard deviation of those expert scores on the test split;
  the proprietary 18-score ensemble's per-dimension Pearson
  correlation on this benchmark from \citet{BenoitEtAl2025},
  Figure~1; and our best Pearson correlation on the same held-out
  documents. Our numbers come from the
  joint multi-dimension run except for economic and
  decentralization, where a single-dimension run gives the better
  number. We meet or beat the ensemble on four of six dimensions;
  social and environment fall slightly short.}
  \label{tab:min-manifesto-data-summary}
\end{table}

\paragraph{Held-out evaluation.}
For each dimension we compute the Pearson correlation between $f$'s
prediction on each manifesto's root summary and the expert-survey
mean for that manifesto, on the test split. The macro is the
unweighted mean across the six dimensions.

\paragraph{Training metric.}
The training metric is per-node agreement between predictions and
a \emph{teacher trace} computed once before training. The teacher
trace is the initial scorer prompt's prediction on the raw text
spanned by each tree node, with no summarization in between, cached
ahead of any optimization. For each training-tree node we then have
a target six-dimension score (from the teacher trace) and a
predicted score (from the trained $f$ on whatever input that node
carries: raw text when training $f$, summary when training $g$).
The metric is the mean across active dimensions of
$\max\!\left(0,\, 1 - \left|\text{predicted} - \text{target}\right|\right)$
with both quantities normalized to $[0, 1]$, and MIPROv2 maximizes
it. In the framework's notation this corresponds to $J_\lambda$
(Eq.~\ref{eq:min-J-lambda}) with the teacher trace standing in for
the oracle at every node. \fixme{Confirm and insert the exact
$\lambda$ used in the runs reported below.}

The cached teacher predictions are what give the summarizer $g$ a
per-node training signal: each internal node has a target score, so
$g$ can be trained to produce summaries that yield the right score
(for a fixed scorer $f$) at every level of the tree, not just at
the document root where the expert mean is available. The
framework's $\lambda$ controls how strongly that per-node signal
weighs against root-level oracle agreement. At $\lambda = 0$ the
local-law term drops out of the loss and $g$'s only training signal
is its downstream effect on $f$'s root prediction; the cached
teacher predictions remain the per-node targets the audit gap below
is measured against, but they no longer enter the training loss. In
that case $g$'s contribution reduces to producing summaries that
make $f$'s root-level evaluation more accurate, with no direct
supervision at internal nodes.

This estimate is a proxy for the framework's $f^*$-based local-law
residual, not the residual itself. The teacher trace is the initial
scorer $f$ applied to raw text, not the true oracle $f^*$, so when
$f$ and $f^*$ disagree their local-law contributions disagree as
well and the trained system optimizes against a distorted target.
The residual against $f^*$ itself is only directly observable when
the oracle is available online: when expert scores can be queried
on demand at every internal node. When node-level oracle access is
unavailable, as here, we estimate the residual against the teacher
trace and read the audit gap below as the visible part of the
proxy's bias.

\paragraph{Audit gap.}
The audit gap is our empirical estimate of the local-law residual
on the test split. Internal Pearson is the correlation between the
trained $f$'s predictions and the teacher-trace targets, taken
across all scored internal nodes; external Pearson is the
correlation at the root only, against expert means. The audit gap
is the difference between the two. A small gap means the scorer's
agreement with the teacher carries to gold; a large gap means the
scorer fits its teacher better than the teacher fits gold.

\fixme{The gap between $f$ and $f^*$ propagates into confidence
intervals on the downstream Pearson correlation, but the gap itself
is estimable on the test split from the conditional average error
between $f$'s root predictions and the expert means, so the
contribution of this distortion is measurable rather than hidden.
Combining the offline teacher-tracked residual with online oracle
queries (when, e.g., a sample of internal nodes can be expert-scored
directly) is then a matter of balancing two estimators of the same
underlying residual; a principled balance is left for future
work.}

\subsection{Single-Dimension Training}\label{sec:min-manifesto-singledim}

We run the setup above first in single-dimension mode: a separate
$(f, g)$ pair for each policy dimension, with that dimension's
expert mean as the only document-level target. Sweeping leaf sizes
from $256$ to $8{,}192$ tokens, we report two points per leaf for
each dimension — the state after round 1 and the best round at
that leaf size — together with three Benoit reference lines:
split-expert agreement, the proprietary 18-score ensemble, and
the best open-weight single-model replication. The substantive reading is that single-dimension
training \emph{recovers the proprietary-LLM ceiling} on each
dimension: a single open-weight model trained against one scorer
at a time matches what an 18-score frontier-LLM ensemble achieves
on the same benchmark.

\paragraph{Economic policy.}
Figure~\ref{fig:v5-manifesto-economic-leafsize} shows the leaf
sweep. External Pearson moves within a roughly $0.04$-wide range
across $32\times$ in leaf size, in line with
Corollary~\ref{cor:schedule}. The 256-token leaf reaches $0.909$
after round 1, $0.029$ above split-expert agreement ($0.880$).
The audit gap stays between $0.056$ and $0.108$ throughout. Round
1 is competitive at every leaf; subsequent rounds only marginally
move the score, so the best round and round 1 sit on top of each
other for most rungs. The best round at every leaf size sits above
Benoit's proprietary ensemble and best open-weight replication; at
the smaller leaves the run also clears split-expert agreement.

\begin{figure}[!htb]
  \centering
  \includegraphics[width=0.72\linewidth]{economic_singledim_leaf_invariance.pdf}
  \caption{\textbf{Single-dimension economic run.} External Pearson
  on the held-out split across leaf sizes, at two points per leaf:
  after round 1 and the best round at each leaf size. Reference
  lines from \citet{BenoitEtAl2025}: split-expert agreement, the
  proprietary 18-score ensemble, and the best open-weight
  single-model replication.}
  \label{fig:v5-manifesto-economic-leafsize}
\end{figure}

\paragraph{Decentralization.}
Decentralization is the dimension where Benoit's existing LLM
systems sit furthest below expert agreement.
Figure~\ref{fig:v5-manifesto-decentralization-leafsize} shows the
leaf sweep. External Pearson ranges from $0.435$ to $0.557$ across
the completed leaves; the $256$-token rung is best at $0.557$ after
round 1, and the $1{,}024$-token rung is next-best at $0.536$ at
its best round. The audit gap ranges from $0.373$ to $0.468$,
considerably larger than the economic run's $0.056$-to-$0.108$
range, marking the dimension as teacher-bounded: even at the best
round the scorer is fitting its teacher trace better than gold.
The best round still clears the proprietary 18-score ensemble
($0.490$) at the $256$- and $1{,}024$-token rungs; the
$512$-token rung's best ($0.476$) sits between the ensemble and
the best open-weight single-model replication ($0.450$).

\begin{figure}[!htb]
  \centering
  \includegraphics[width=0.72\linewidth]{decentralization_singledim_leaf_invariance.pdf}
  \caption{\textbf{Single-dimension decentralization run.} External
  Pearson correlation on the held-out split across leaf sizes, at
  two points per leaf: after round 1 and the best round at each
  leaf size. Reference lines from \citet{BenoitEtAl2025}:
  split-expert agreement, the proprietary 18-score ensemble, and
  the best open-weight single-model replication.}
  \label{fig:v5-manifesto-decentralization-leafsize}
\end{figure}

\subsection{Combined Multi-Dimension Training}\label{sec:min-manifesto-combined}

We now turn to joint training, the universal-summary regime
introduced earlier: one $g$ produces a single summary per
manifesto, and one scorer returns six per-dimension scores from
that summary. $g$ is trained against the average of the six
per-dimension $f$-rewards; $f$ is trained dimension-aware. The
leaf sweep covers all six rungs from $256$ to $8{,}192$ tokens.
Figure~\ref{fig:v5-manifesto-combined-leafsize} plots external
Pearson correlation per dimension across the sweep, with round 1
and the best round at each leaf size, and three Benoit reference
lines per panel. Table~\ref{tab:min-manifesto-audit-read} reports
the best-round result per dimension.

\begin{figure}[!htb]
  \centering
  \includegraphics[width=\linewidth]{manifesto_fg_combined_ladder_f1g0_f1g1.pdf}
  \caption{\textbf{Joint multi-dimension run, per-dimension
  trajectories.} External Pearson on the held-out split across leaf
  sizes for each dimension, sliced from the same joint $(f, g)$
  training. Two series per panel: round 1 and the best round at
  each leaf size. Reference lines per panel from
  \citet{BenoitEtAl2025}: split-expert agreement, the proprietary
  18-score ensemble, and the best open-weight single-model
  replication.}
  \label{fig:v5-manifesto-combined-leafsize}
\end{figure}

\begin{table}[!htb]
  \centering
  \begin{tabular}{lccc}
    \toprule
    Dimension        & After round 1 & Best round & Audit gap \\
    \midrule
    EU integration   & $0.965$ & $0.965$ & $0.013$ \\
    Immigration      & $0.934$ & $0.938$ & $0.033$ \\
    Economic         & $0.876$ & $0.886$ & $0.088$ \\
    Social           & $0.884$ & $0.886$ & $0.081$ \\
    Environment      & $0.795$ & $0.795$ & $0.085$ \\
    Decentralization & $0.461$ & $0.461$ & $0.326$ \\
    \bottomrule
  \end{tabular}
  \caption{\textbf{Per-dimension peaks of the joint multi-dimension
  run.} Each row reports the highest external Pearson on the
  held-out split for that dimension after one alternation round and
  across all completed rounds (best across all leaf sizes in each
  case). The audit gap (internal minus external Pearson) is
  reported at the best-round result.}
  \label{tab:min-manifesto-audit-read}
\end{table}

Three of the six dimensions clear the proprietary 18-score
ensemble in the joint run: EU integration and immigration both also
clear split-expert agreement, and economic policy clears the
ensemble by a small margin. Social and environment fall slightly
short of the ensemble, by roughly $0.025$ each. Decentralization
falls furthest short, at $0.461$.

Where a single-dimension run exists for comparison, the picture is
sharper. Economic's joint peak ($0.886$ at leaf $4{,}096$) sits
within $0.002$ of the single-dimension peak ($0.888$ at leaf
$2{,}048$); the universal summary captures essentially the same
signal here as a dedicated economic summary. Decentralization's
joint peak ($0.461$ at leaf $8{,}192$) sits $0.096$ below the
single-dimension peak ($0.557$ at leaf $256$); the single-dimension
peak also clears the proprietary ensemble on this dimension by
$0.067$.

On decentralization the joint run shows the same shape at every
leaf size. External Pearson stays between $0.30$ and $0.39$, the
audit gap stays around $0.45$, and the stage-by-stage trajectory
oscillates without converging: a $g$-update lifts the score and
the next $f$-update pulls it back. Smaller leaves do not narrow
the gap and bigger leaves do not either. The mechanism is
multi-task interference. Averaging the six $f$-rewards gives any
single dimension a fraction of the gradient when its target points
somewhere different from the others, so any $g$-step that picks up
that dimension's signal gets unwound by the next $f$-step
regardless of leaf size. Removing the averaging (the
single-dimension run, \S\ref{sec:min-manifesto-singledim})
lifts external Pearson from $0.461$ to $0.557$.

The audit gap on the joint run flagged the breakdown without
needing the single-dimension result. On the other five dimensions
the joint audit gap stays below $0.09$; decentralization's $0.326$
is more than three times that. Joint training is the right default
whenever every dimension's audit gap stays in the range the
well-behaved dimensions show; when one dimension blows past that
range, fall back to a single-dimension run for that dimension.

\subsection{Discussion}\label{sec:min-manifesto-takeaways}

The single-dimension and joint runs above land two of the
framework's predictions side by side, and it is worth pausing on
each before moving on.

The first is schedule invariance (Corollary~\ref{cor:schedule}),
which holds across the dimensions we have measured. Best external
Pearson per dimension moves by between $0.014$ and $0.048$ on the
joint run across $32\times$ in leaf size, and by less than $0.04$
on the single-dimension economic run across the same range. Even
decentralization, where the joint run leaves a large audit gap,
moves by less than that gap across the leaf sweep. The chunk
schedule is therefore a free parameter for the practitioner: pick
whatever leaf size makes the audit convenient, and the realized
measurement does not move.

The second is that the audit gap is the operational handle. The
audit gap measures, on the test split, how far the trained
scorer's agreement with the teacher trace exceeds its agreement
with gold. On decentralization the joint run had a gap of $0.326$,
more than three times the upper end of the range the other five
dimensions show. The single-dimension run on the same dimension
kept the gap at a similar magnitude (between $0.373$ and $0.468$
across leaves) but lifted external Pearson from $0.461$ to $0.557$
at its best leaf, enough to clear the proprietary ensemble. The gap does not shrink to zero in the
single-dimension run because decentralization remains
teacher-bounded: the teacher trace itself encodes a noisier
mapping from text to gold here than on the other dimensions. The
framework's diagnostic separates that ceiling from the
joint-training failure cleanly, and Section~\ref{sec:min-audit}
turns the same internal estimates into a deployable drift bound on
the document-level score.
